\section{水素原子モデル}

\subsection{トムソンの原子模型とラザフォードの原子模型}
物質は原子という，それ以上分割できない粒子から構成されているという概念は古くから存在したが，1897年にトムソンによって電子の比電荷が測定され，同年にゼーマンによって電子が原子内に含まれていることが確認されると，原子内にどのような電子が存在するかに興味が持たれた．
現在ではすでにこの問題については結論が得られているが，当時は次のようなモデルが考えられていた．

電子の質量は原子に比べて非常に小さく，また原子全体は電気的に中性であることから，原子から電子を取り除いた残りの部分とは，原子の大部分の質量を持ち，かつ電子の負電荷を打ち消す正電荷を持ったものだと考えられる．
また，原子はある特定の波長・振動数の電磁波を吸収したり放出したりすることが当時知られていた．

\begin{wrapfigure}[6]{r}[5mm]{60mm}
  \centering
  \vspace{-10mm}
  \includegraphics[keepaspectratio,width=60mm]
  {text2_p102_fig1.png}
\end{wrapfigure}

そこで，トムソンは，正電荷が球状の原子いっぱいに雲のように一様に分布し，電子がこの中に埋め込まれて安定な状態でつり合って静止しているモデル（ブドウパンモデルとも呼ばれる）を考えた．
この場合，外部からの何かしらの作用によって電子が安定位置からずれると，つり合いの位置を中心として振動し，その振動数と一致した電磁波を吸収・放出すると考えられる．

当時は直接的に原子を観察する手段など存在しなかったが，ラザフォードは次のような実験からトムソンのモデルでは実験事実を説明できないことを示し，上の図にあるようなラザフォードのモデルを提案した．

\begin{wrapfigure}[15]{r}[5mm]{60mm}
  \centering
  \vspace{-5mm}
  \includegraphics[keepaspectratio,width=60mm]
  {text2_p102_fig2.png}
\end{wrapfigure}

ラザフォードの弟子であるガイガーとマースデンは，放射性元素ラジウムから放出される$\alpha$粒子（=ヘリウムイオン$ \rm ^4_2He^{2+}$）を厚さ約$0.5[{ \mu \rm m}]$の金箔に垂直に当て，その後の$\alpha$粒子の進行方向について調べた．
すると，大半の$\alpha$粒子はほとんど直進し，$10^5$個に1個程度の割合でごく一部の$\alpha$粒子が非常に大きく曲げられることがわかった．

トムソンのモデルでは，原子の内部でも電荷がほぼ一様に分布していることから，$\alpha$粒子を入射しても，大きな反発力を受けず，直進するか，曲がるとしても，非常に質量の軽い電子からの反発しよるものであるため，曲がり方はわずかになるはずである．
進行方向に対して$90^\circ$以上方向を変えるような大きな散乱をするためには，$\alpha$粒子が大きな反発力を受ける必要がある．

そこで，図にあるようなラザフォードのモデルを想定すると，大半の$\alpha$粒子は直進するが，ごくわずかの$\alpha$粒子が，極めて狭い部分に集中した正電荷のすぐそばを通り，大きな斥力を受けて進路を大きく変えたと考えられ，実験結果をうまく説明することができる．

以上から，ラザフォードは以下のような結論に達した．

\begin{description}
\item[\MARU{1}] 原子の中心に，原子のほぼ全質量を持つ原子核が存在する．
\item[\MARU{2}] 原子番号$Z$の原子核は$+Ze$[C]の正電荷を持っている．
また，原子核の半径は電子の円運動半径の約$1/10000$である．
\item[\MARU{3}] 原子核の周囲を$Z$個の電子が回っており，原子核の正電荷を中和して電気的に中性になっている．
\end{description}

ところが，一見正しそうに見えるこのラザフォードのモデルにも次のような難点があった．
先にも述べたように，原子は特定の波長の電磁波を吸収・放射（輻射）することが知られている．
例えば水素原子の気体は下図のように，紫外線領域，可視光領域，赤外線領域にそれぞれ群をなして，ある特定の波長の電磁波のみを放射する．

\begin{figure}[h]
\begin{center}
\includegraphics[width=120mm]{text2_p103_fig1.png}
\end{center}
\end{figure}

表にあるように，紫外部に現れる一群をライマン系列，可視部に現れる一群をバルマー系列，赤外部に現れる一群をパッシェン系列という．

\begin{wrapfigure}[6]{r}[5mm]{60mm}
  \centering
  \vspace{-10mm}
  \includegraphics[keepaspectratio,width=40mm]
  {text2_p103_fig2.png}
\end{wrapfigure}

ラザフォードのモデルを水素原子に適用してみよう．
水素原子は$+e$[C]の電荷を持つ原子核のまわりを$-e$[C]の電荷を持つ電子がクーロン力を受けて等速円運動していると考えられる．
このとき，電子の速さを$v$[m/s]，質量を$m$[kg]，軌道半径を$r$[m]とすると運動方程式は，
\[m\dfrac{v^2}{r}=k\dfrac{e^2}{r^2}\]
であり，電子の持つエネルギーは運動エネルギーとクーロン力による位置エネルギーの2つの和であるから，
\[E=\dfrac{1}{2}mv^2-k\dfrac{e^2}{r}=-\dfrac{ke^2}{2r} \cdots (☆)\]
となる．

\begin{wrapfigure}[10]{r}[5mm]{60mm}
  \centering
  \vspace{-5mm}
  \includegraphics[keepaspectratio,width=50mm]
  {text2_p104_fig1.png}
\end{wrapfigure}

電磁気学によると，電荷を持ったものが円運動すると，その円運動の回転数に等しい振動数を持つ電磁波を前方に放出して，エネルギーを失うことが知られている．
$E$が減少すると，(☆)式から$r$が減少し，電子はやがて核にくっついてしまい，トムソンの原子模型と同じことになってしまう．
すなわち，{\bf ラザフォードモデルでは原子半径が定まらない}という難点がある．

また，発せられる電磁波の振動数は，
\[\nu=\dfrac{v}{2\pi r}\varpropto \dfrac{1}{\sqrt{r^3}}\]
となるため，原子半径$r$の減少と共に振動数$\nu$は連続的に大きくなる．
しかし実際には，ある特定の波長（振動数）の電磁波しか出さない．
すなわち，{\bf ラザフォードモデルでは原子の発する電磁波が連続的に変化してしまう}という難点がある．

これらのことより，ラザフォードのモデルでは理論的には原子が安定に存在できないことになってしまう．
このような問題が生じてしまうのは，通常の電磁気学を原子内の電子という非常に微細な粒子に対して適用してしまったからであり，次に述べるボーアの理論を適用することによりこれらの難点は全て解決できる．

\newpage
\subsection{ボーアの水素原子モデル}
\begin{wrapfigure}[20]{r}[5mm]{60mm}
  \centering
  \vspace{-10mm}
  \includegraphics[keepaspectratio,width=60mm]
  {text2_p104_fig2.png}
    \includegraphics[keepaspectratio,width=65mm]
  {text2_p105_fig1.png}
\end{wrapfigure}



先に述べたラザフォードの難点を，ボーアは次の3つの仮定を持ち出すことで見事に解決することに成功した．
\begin{description}
\item[\MARU{1}] {\bf 電子が原子内で持つエネルギーは連続的に全ての値を取りうるのではなく，いくつかの特定のエネルギー値しか持たない．また，そのエネルギーを持った状態にある限り，電磁波の放射は起こらない．}
この状態のことを{\bf 定常状態（安定状態）}といい，この安定状態のエネルギー値のことを{\bf エネルギー準位}という．
\item[\MARU{2}] {\bf 電子の軌道円上に電子波の定常波が実現するときに安定状態となりうる．}

すなわち，電子のド・ブロイ波長を$\lambda$[m]とするとき，
\[2\pi r=n\lambda \hspace{4mm}(n\in \mathbb{N})\]
が成立するとき電子は安定状態となり，この条件を満たすような半径$r$のところにしか電子は安定に存在できない．
これを{\bf 量子条件}という．
また，$n$は{\bf 量子数}といい，この$n$の値に応じて$r$，$E$が決まるので，これらを以後それぞれ$r_n$，$E_n$と表すことにする．
\item[\MARU{3}] {\bf 電子がエネルギー準位}\mbox{\boldmath$E_n$}{\bf の定常状態から}\mbox{\boldmath$E_{n^\prime}$}{\bf の定常状態に移るときのエネルギー差}\mbox{\boldmath$|E_n-E_{n^\prime}|$}{\bf に等しいエネルギーの光子のみ吸収または放出することができる．}

すなわち，$|E_n-E_{n^\prime}|=h\nu$を満たす振動数$\nu$の光子を吸収または放出する．
これを振動数条件という．
このため，吸収・放出される光の振動数は限定される．



右図のように，電子がエネルギー準位間を移動するのに伴って光子が吸収・放出される．
光子を吸収して電子がよりエネルギーの低い準位に移動することを{\bf 放射}（もしくは{\bf 輻射（ふくしゃ）}）という．
また，このように電子がエネルギー準位間を移動することを{\bf 遷移（せんい）}という．
\end{description}

この3つの仮説を立てることにより得られる，水素原子の電子の軌道半径やエネルギー準位から計算すると，水素原子の放出する電磁波を見事に説明することができる．

最も簡単な水素原子のモデルとして有名なものにボーアの模型がある．
このボーアの模型は，電荷$+e{{\rm C}}$の水素原子核の周りを電荷$-e$[C]，質量$m$[kg]の電子が等速円運動をしているものとして与えられる．
このモデルで運動方程式$m\dfrac{v^2}{r}=k\dfrac{e^2}{r^2}$，量子条件$2\pi r=n\lambda$，エネルギー$E=\dfrac{1}{2}mv^2-k\dfrac{e^2}{r}$の3つの式を組み合わせて計算すると，水素原子の電子の軌道半径$r_n$およびエネルギー準位$E_n$は，次のようになる．
$$r_n=\dfrac{h^2}{4\pi^2kme^2}\dot n^2\hspace{1mm}[{\rm m}] \hspace{5mm} e_n = - \dfrac{2 \pi^2 k^2 m e^4}{h^2}・ \dfrac{1}{n^2} \hspace{1mm} [{\rm J}]$$
ここで，$n$は自然数であるから飛び飛びの値しか取ることができず，また，
$$r_1<r_2<r_3<\cdots ，\hspace{5mm} E_1<E_2<E_3<\cdots $$
なる関係になっている．
一般にエネルギーは低いほど安定であるから，エネルギー準位$E_1$に電子がいるとき，すなわち$n=1$ の軌道に電子がいるとき，水素原子は最も安定である．
これを基底状態という．
それに対して，$n\geqq 2$の軌道に電子がいるときのことを励起状態という．

\begin{wrapfigure}[10]{r}[5mm]{60mm}
  \centering
  \vspace{-5mm}
  \includegraphics[keepaspectratio,width=50mm]
  {text2_p106_fig1.png}
\end{wrapfigure}
基底状態における水素の原子半径$r_1$を{\bf ボーア半径}$a_0$といい，これを計算すると，
$$a_0=\dfrac{h^2}{4\pi^2kme^2} \fallingdotseq 5.3\times 10^{-2} [{\rm nm}]$$
となり，水素原子の大きさは，$0.1$[nm]程度であるとわかる．


また，このときのエネルギー$E_1$は
$E_1=-\dfrac{2\pi^2k^2me^4}{h^2} \fallingdotseq -13.6 [{\rm eV}]$
である．
よって，これらを用いると，
$$r_n=n^2a_0,\hspace{5mm} E_n=-\dfrac{13.6}{n^2}[eV]$$
と書くことができる．
そのため，水素原子のエネルギー準位は上に行くほど間隔が狭まっていく．

また，$n\to \infty$のときは$r_\infty =\infty$，すなわち電子がはるかかなたへと去った場合，すなわちイオン化した場合に対応するので，$E_\infty -E_1\fallingdotseq 13.6[{\rm eV}]$は水素原子のイオン化エネルギーである．

\begin{matome}{水素原子のエネルギー準位}
 \hspace{2mm} 水素原子のエネルギー準位$E_\infty -E_1\fallingdotseq 13.6[{\rm eV}]$であり，
 $E_n=-\dfrac{13.6}{n^2}[{\rm eV}]$
 \end{matome}
 
  \newpage

\begin{figure}[h]
\begin{center}
\includegraphics[width=170mm]{prob2_p61.png}
\vspace{150mm}
\end{center}
\end{figure}

\newpage

\begin{figure}[h]
\begin{center}
\includegraphics[width=170mm]{prob2_p62.png}
\vspace{150mm}
\end{center}
\end{figure}

\newpage

\begin{figure}[h]
\begin{center}
\includegraphics[width=170mm]{prob2_p180_2.png}
\vspace{150mm}
\end{center}
\end{figure}

\newpage

\begin{figure}[h]
\begin{center}
\includegraphics[width=170mm]{prob2_p181.png}
\vspace{150mm}
\end{center}
\end{figure}

\newpage

\begin{figure}[h]
\begin{center}
\includegraphics[width=170mm]{prob2_p182_1.png}
\vspace{150mm}
\end{center}
\end{figure}
